% ECONOMICS_PROBLEM: {"id": "I18-263", "title": "Antibiotic innovation incentives", "jel_code": "I18", "source_id": "OP-0446"}
% !TeX program = xelatex
\documentclass[11pt,a4paper]{article}
\usepackage[margin=23mm]{geometry}
\usepackage{fontspec}
\setmainfont{Latin Modern Roman}
\usepackage{amsmath,amssymb,mathtools}
\usepackage{xcolor,enumitem,booktabs,longtable,array,xurl,hyperref}
\definecolor{muted}{HTML}{53636C}
\hypersetup{colorlinks=true,urlcolor=blue,linkcolor=blue}
\setlength{\parindent}{0pt}
\setlength{\parskip}{5pt}
\newcommand{\R}{\mathbb R}
\newcommand{\E}{\mathbb E}
\newcommand{\N}{\mathbb N}
\newcommand{\ind}{\mathbf 1}
\newcommand{\Var}{\operatorname{Var}}
\newcommand{\Cov}{\operatorname{Cov}}
\newcommand{\supp}{\operatorname{supp}}
\DeclareMathOperator*{\argmin}{arg\,min}
\DeclareMathOperator*{\argmax}{arg\,max}
\newcommand{\entrymeta}[3]{{\sffamily\small\color{muted}\textbf{\texttt{#1}}\quad|\quad #2\par #3\par}}
\newcommand{\Problemref}[1]{\texttt{#1}}
\newcommand{\JELref}[2]{\texttt{#2} (JEL \texttt{#1})}
\begin{document}
\section*{Antibiotic innovation incentives}
\textbf{Problem ID:} \texttt{I18-263}\quad \textbf{JEL:} \texttt{I18}\par
% BEGIN ECONOMICS STATEMENT
\label{problem:OP-0446}
{\sffamily\small\color{muted}Original subject group: Health economics\par}
\label{definitions:problem:30007728}
\textbf{Audit classification:} Background; proposed extension. Prevented duplicate valuation of resistance health losses and separated payments from resource costs.
\subsection*{Definitions and precise research target}
\textit{Editorial formalization.} Consider developers of antibiotics for a prespecified set of resistant pathogens in European health systems. A policy \(d\in\mathcal D\) is a fully specified combination of research grants, milestone prizes, subscription payments, and stewardship requirements; \(\mathcal D\) is a finite feasible set. Let \(A(d)\) be the number of qualifying antibiotics approved within ten years, \(H(d)\) discounted patient health gained, \(R(d)\) resistance burden after usage responses, and \(K(d)\) discounted real public and private resource cost. For positive valuations \(\lambda,\mu\), define
\[d^*\in\arg\max_{d\in\mathcal D}E[\lambda H(d)-\mu R(d)-K(d)]\quad\text{subject to}\quad E[A(d)]\geq a_0,\]
where \(a_0\) is a stated minimum innovation target. Report infeasibility if no contract reaches the innovation target, and require finite expected objective terms. Model entry, scientific failure, information available to the payer, and manufacturer incentives under each contract. Specify which quantities can be estimated from existing programs and which require bounds or elicited scientific probabilities. Include cross-border resistance externalities and distinguish payments from resource costs. Ask whether revenue independent of sales can produce valuable innovation while preserving low appropriate usage. This is a proposed mechanism and quantitative design problem; the original policy review was not verified as explicitly posing this exact optimization problem.

A qualifying antibiotic has prespecified pathogen coverage, novelty, clinical effectiveness, and approval criteria. Define \(A(d)\) as its approval count, \(H(d)\) as quality-adjusted patient-years gained, and \(R(d)\) as an additional resistance externality valued in declared units. If resistance already reduces \(H\), do not value the same loss again in \(R\): either define \(R\) as a separate ecological/nonhealth damage or include all resistance health consequences in \(H\) and omit \(\mu R\). \(K\) is real R\&D, delivery and administration cost; grants and subscriptions are budget transfers and require a separate payer budget constraint. The menu includes contract terms and stewardship enforcement, not just names of instruments. With finite integrable feasible choices the optimizer exists. The substantive editorial task is estimating innovation, use and resistance responses under these contracts. All expectations use a stated baseline population and probability law. Identification means every admissible data-generating model with the same observed-data law yields the same displayed target; otherwise report the identified set. Exact equations, horizons and assumptions are editorial, rather than source-given canonical conjectures.
\subsection*{Variants and assumption changes}
\begin{enumerate}
\item \textbf{Editorial sales-linked variant.} Compare per-dose reimbursement and subscription at equal expected fiscal expenditure. Estimate use and resistance changes; equal budgets do not imply equal resource costs.
\item \textbf{Editorial international variant.} Include low-income-country access and global resistance propagation. Compare coordinated and unilateral funding under stated spillover rules.
\end{enumerate}
\subsection*{References and source audit}
\textbf{1.} Review compares incentive options; exact innovation-constrained optimizer is editorial. Locator: Lancet Regional Health–Europe 33, 100705; Summary, p.1; author-repository PDF. PMC follow-ups met a challenge page; identical primary author-repository article inspected. URL: \url{https://pmc.ncbi.nlm.nih.gov/articles/PMC10403717/}.
\begin{enumerate}\raggedright
\item\hypertarget{definitions:source:30007728:1}{} Michael Anderson; Dimitra Panteli; Robin van Kessel; Gunnar Ljungqvist; Francesca Colombo; Elias Mossialos. 2023. \emph{Challenges and opportunities for incentivising antibiotic research and development in Europe}. Lancet Regional Health–Europe 33, 100705; Summary, p.1; author-repository PDF.
\url{https://pmc.ncbi.nlm.nih.gov/articles/PMC10403717/}.
DOI: \url{https://doi.org/10.1016/j.lanepe.2023.100705}.
\end{enumerate}

\subsection*{Common conventions: Source mathematical conventions}

These conventions apply to each entry unless explicitly replaced.

\textbf{Models and identification.} A primitive model \(m\) specifies all probability laws, feasible choices, information, equilibrium and measurement rules used by its target. The admissible class \(\mathcal M\) is fixed independently of outcomes. If \(P_O(m)\) is its observable probability law and \(\tau(m)\) the target, its identified set at observable law \(P\) is
\[
 \mathcal I_\tau(P)=\{\tau(m):m\in\mathcal M,\ P_O(m)=P\}.
\]
Point identification means that this set is a singleton. An empty set signals incompatibility of data and assumptions. Coordinate bounds need not describe the joint set. Bounds are sharp only if every retained target is attainable or arbitrarily approachable by compatible models. Finite bounds require explicit restrictions on unobserved quantities; unrestricted confounding, hidden wealth, extrapolation or arbitrary equilibrium selection can yield uninformative sets.

\textbf{Causal interventions.} A potential outcome \(Y_i(p)\) is the outcome under a complete policy regime \(p\), including timing, financing, information and market spillovers. The target population and units are fixed before treatment. With interference, use the complete assignment vector or a justified exposure mapping; individual no-interference notation alone cannot identify an economy-wide intervention. A difference of mean potential outcomes does not require the cross-world joint distribution. Conditioning on post-treatment migration, survival, enrollment or employment changes the estimand and needs separate justification.

\textbf{Statistical statements.} An estimator, confidence set, forecast or test specifies a sampling law, model class and loss/coverage criterion. Uniform \((1-\alpha)\)-coverage means \(\inf_{m\in\mathcal M}\Pr_m\{\tau(m)\in C_n\}\geq1-\alpha\), or an explicitly asymptotic version. The whole parameter space trivially has coverage, so informativeness requires a stated width, risk or power objective. A forecasting improvement is relative to a named benchmark, information set and evaluation population.

\textbf{Equilibrium and optimization.} An equilibrium comprises feasible plans, optimality, beliefs where needed, budget constraints and all market/resource clearing. The primitives or a stated benchmark define each correspondence \(\mathcal E(p;m)\). Nonemptiness is assumed or proved, never inferred just from compact policy sets. A selected-equilibrium target uses a declared selection rule; a robust target quantifies over all permitted equilibria. Use suprema or infima unless attainment is established. An \(\varepsilon\)-optimal policy is feasible and within \(\varepsilon\) of the finite optimum value. Report an empty feasible set separately.

\textbf{Welfare and accounting.} Normative utility, interpersonal normalization, social weights, numeraire, discounting and the use of public receipts are primitives. Behavioral utility need not coincide with the welfare criterion. Household welfare includes ownership income and rents once; adding profits again to compensating variation generally double-counts them. A monetary equivalent is defined by an explicit transfer and price convention, with monotonicity and range conditions. Average effects, mechanism contributions, transfers and welfare losses are different targets. Infinite sums require convergence and dynamic feasible plans require terminal or no-Ponzi conditions.

\textbf{Source versions.} A working-paper date, revision date and journal publication year are distinguished. A DOI for a journal version is not automatically the DOI of an earlier manuscript. Locators refer to the stated version and printed pages where specified. A metadata-only check does not verify a claim about a full-text conclusion. Every entry's source subsection narrows the provenance claim accordingly.
% END ECONOMICS STATEMENT
\end{document}
